% !TeX root = ../tango.tex
\section{Related Work}
\label{sec:related-work}

\begin{table}[t]
  \renewcommand{\arraystretch}{1}
  \setlength{\tabcolsep}{1.5pt}
  \caption{Comparison of \tango{} with representative serving systems.}
  \label{tab:related-work-comparison}
  \centering
  \scriptsize
  \resizebox{\columnwidth}{!}{%
  \begin{tabular}{|c|c|c|c|c|}
    \hline
    \bfseries & \textbf{\makecell{Continuous\\Batching}} & \textbf{\makecell{Heterogeneous\\SLOs}} & \textbf{\makecell{Runtime\\Reprioritization}} & \textbf{\makecell{Predictive Token\\Budgeting}} \\
    \hline
    \bfseries vLLM & \cellcolor{FadedFlora}\CheckmarkBold & & & \\
    \hline
    \bfseries SOLA & \cellcolor{FadedFlora}\CheckmarkBold & & \cellcolor{FadedFlora}\CheckmarkBold & \\
    \hline
    \bfseries JITServe & \cellcolor{FadedFlora}\CheckmarkBold & \cellcolor{FadedFlora}\CheckmarkBold & \cellcolor{FadedFlora}\CheckmarkBold & \\
    \hline
    \bfseries \tango{} & \cellcolor{FadedFlora}\CheckmarkBold & \cellcolor{FadedFlora}\CheckmarkBold & \cellcolor{FadedFlora}\CheckmarkBold & \cellcolor{FadedFlora}\CheckmarkBold \\
    \hline
  \end{tabular}
  }
  \vspace{-3mm}
\end{table}

\textbf{LLM Serving Systems.}
Recent studies have developed efficient LLM serving systems that improve inference throughput and latency through scheduling and memory optimizations.
Orca introduced iteration-level scheduling, allowing requests to enter and leave a batch between decoding iterations to improve both utilization and responsiveness.
vLLM further improved serving throughput through PagedAttention, which reduces KV-cache fragmentation and enables flexible memory sharing across requests.
SGLang exploits prefix sharing and structured execution to accelerate multi-call LLM workloads.
These systems establish continuous batching and efficient memory management as the foundation of modern LLM serving.
However, throughput-oriented schedulers commonly use an SLO-agnostic request order and a fixed upper bound on the number of tokens processed in each iteration, without adapting either decision to the heterogeneous TTFT and TPOT requirements of individual requests.


\textbf{Prefill--Decode Scheduling.}
The prefill and decode phases exhibit distinct execution characteristics: prefill processes many prompt tokens in parallel, whereas decode repeatedly processes one token per request.
Mixing the two phases improves GPU utilization but can prolong an iteration and delay the next tokens of ongoing requests.
Sarathi-Serve mitigates this interference through chunked prefill and stall-free batching, using decode-maximal batches to improve serving capacity under tail-latency constraints.
DistServe instead places prefill and decode on different GPUs and separately optimizes their resource allocation and parallelism under TTFT and TPOT constraints.
These approaches regulate interference through configured chunk sizes, phase-level batching policies, or resource separation.
They do not directly translate the remaining latency slack of active requests into an iteration-specific token budget for a colocated mixed batch.


\textbf{SLO-Aware LLM Serving.}
Recent systems explicitly optimize the fraction of requests satisfying latency SLOs rather than raw throughput alone.
SOLA formulates a fine-grained scheduling space that adjusts both request execution order and iteration workload according to request and system states, thereby balancing TTFT and TPOT.
JITServe estimates uncertain request characteristics and allocates each request sufficient serving bandwidth while grouping compatible requests to improve batch efficiency and goodput.
However, existing systems do not adjust the per-step token allocation within a mixed batch according to the request-specific TTFT and TPOT slack.
Table~\ref{tab:related-work-comparison} compares \tango with representative serving systems.
